\subsection{27.01.2026 (Freistunde)}

\subsubsection{Ziel}

Ziel des Tages war es, die Taupunktberechnung für Python
anzupassen und in das bestehende Projekt zu integrieren.

\subsubsection{Durchgeführte Arbeiten}

Zunächst wurden Klassen- und Dateinamen an gängige
Namenskonventionen angepasst, um die Struktur und
Lesbarkeit des Projekts zu verbessern.

Anschließend wurden die Funktionen zur Berechnung des
Taupunkts in Python implementiert und in das bestehende
Projekt integriert.

Die Funktionalität wurde dabei auf die beiden Dateien
\texttt{scann\_berechnen.py} und
\texttt{taupunkt\_berechnen.py} aufgeteilt,
um den Quellcode übersichtlicher zu strukturieren.

In \texttt{scann\_berechnen.py} wurde die Klasse
\texttt{TaupunktLogik} erstellt.
Diese importiert sowohl den DHT11-Scanner als auch die
Taupunktberechnung, verarbeitet die Sensordaten und
übergibt diese an die Berechnungsklasse.

Die berechneten Ergebnisse werden anschließend wieder
an die Hauptlogik zurückgegeben.

Zusätzlich wurde in
\texttt{taupunkt\_berechnen.py} die Klasse
\texttt{TaupunktBerechnung} erstellt.
Diese erhält die Sensordaten als Übergabeparameter und
berechnet daraus die Taupunkttemperatur der Messpunkte.

Zum aktuellen Zeitpunkt wurde die Berechnung jedoch erst
für einen Sensor umgesetzt.

Darüber hinaus wurde die Klasse
\texttt{TaupunktLogik} bereits in
\texttt{taupunkt\_main.py} importiert.
Eine aktive Nutzung erfolgt jedoch noch nicht, da der
zweite Sensor bislang nicht implementiert wurde.

Es wurden keine Tests durchgeführt, da der Joy-Pi aktuell nicht zur Verfügung steht.

\subsubsection{Erkenntnisse}

Es wurde festgestellt, dass Python-Dateien üblicherweise
im Snake-Case benannt werden, während für Klassennamen
die CamelCase-Schreibweise verwendet wird.

Zusätzlich wurde deutlich, dass eine Aufteilung des
Programms in mehrere spezialisierte Klassen die
Wartbarkeit und Übersichtlichkeit des Projekts verbessert.

\subsubsection{Nächste Schritte}

In der nächsten Woche soll der zweite DHT11-Sensor
angeschlossen und in die Taupunktberechnung integriert
werden.

Außerdem sollen die Berechnungsergebnisse in die
Hauptlogik übernommen und anschließend auf dem
LCD-Display ausgegeben werden.